\documentclass[12pt]{article}
\usepackage{graphicx} % Required for inserting images
\usepackage[margin=0.5in]{geometry}
\usepackage{amsmath, amssymb}
\usepackage{mathrsfs}


\usepackage{soul}


% Dr. David Brown Commands
\newcommand{\ds}{\displaystyle}
\newcommand{\R}{\mathbb{R}}
\newcommand{\N}{\mathbb{N}}
\newcommand{\Q}{\mathbb{Q}}
\newcommand{\C}{\mathbb{C}}


% Commands to get script r
\usepackage{calligra}
\DeclareMathAlphabet{\mathcalligra}{T1}{calligra}{m}{n}
\DeclareFontShape{T1}{calligra}{m}{n}{<->s*[2.2]callig15}{}
\newcommand{\scripty}[1]{\ensuremath{\mathcalligra{#1}}}

% Unit commands
\newcommand{\degree}{\text{°}}
\newcommand{\unitSpeed}{\frac{\text{m}}{\text{s}}}
\newcommand{\unitAcceleration}{\frac{\text{m}}{\text{s}^2}}

% Constant commands
\newcommand{\avogadro}{6.022\times10^{23}}

\newcommand{\boltzmannConstK}{1.381\times10^{-23}\frac{\text{J}}{\text{K}}}
\newcommand{\planckConst}{6.626\times10^{-34}\text{J}\cdot\text{s}}

\newcommand{\atmP}{1.01325\times10^5\,\text{Pa}}

% Commands to easily write partial derivatives
\newcommand{\partDeriv}[2]{\frac{\partial #1}{\partial #2}}
\newcommand{\doublePartDeriv}[2]{\frac{\partial^2 #1}{\partial^2 #2}}


\title{Problem 75}
\author{Gabriel Decker}


\begin{document}



\maketitle
\begin{flushleft}



We're given the one-dimensional force $F=F_0\sin(\omega t)$.
By Newton's second law ($\vec{F}=m\vec{a})$, this implies that
$$F_0\sin(\omega t)=ma$$
$$\implies a=\frac{F_0}{m}\sin(\omega t)$$\\~\\

Jerk is the derivative of acceleration with respect to time.
Therefore the jerk $J$ is
$$J(t)=\frac{da}{dt}=\frac{F_0\omega}{m}\cos(\omega t)$$
$J$ is maximum when the cosine function returns its maximum value, 1.
Therefore the maximum jerk is $\ds\frac{F_0\omega}{m}$.




\end{flushleft}

\end{document}
